\documentclass[main.tex]{subfiles}
%subfile :prelim_2.tex
\begin{document}
	\section{Induded Covers}  
    \label{sec:prelim_induced_covers}


{\definition{}
Let $\mathcal{B}=\{B_\alpha: \alpha\in\Delta\}$ be a locally finite
family of subsets of $X$. 
Let $\mathcal{B}^* = \{U_s: s\in [\Delta]^{<\omega}  \text{ and} $
$U_s = X - \bigcup_{\gamma \notin s} \overline{B_\gamma}\} $. 
Then $\mathcal{B}^*$ 
is an open cover of $X$. We call $\mathcal{B}^*$ an 
\uwave{induced cover}\index{induced cover, \S\ref{sec:prelim_induced_covers}} 
of $X$.
}%end of definition

Remark. 
It is easy to see that 
$\mathcal{A} = \{U_s\bigcap \cap_{\alpha\in s} \overline{B_\alpha}: 
s\in [\Delta]^{<\omega}  \text{ and} $
$U_s = X - \bigcup_{\gamma \notin s} \overline{B_\gamma}\}$ 
is a locally finite refinement of $\mathcal{B}^*$. Later, 
we will see that
$\mathcal{B}^*$ is an 
interior-preserving (Definition \ref*{def_prelim3_cp_ip}) directed open cover of $X$. 

{\definition{}
Let $\mathcal{D} = \{D_\alpha : \alpha \in \Delta\}$ be a discrete family of 
closed subsets of $X$. Let $\mathcal{U}_\alpha = X - \bigcup_{\gamma \neq \alpha} D_\gamma$ and let $\mathcal{B}^{*F} = \{\mathcal{V}_S : S \in \mathcal{F}\}$ where $\mathcal{V}_S = \bigcup_{\alpha \in S} \mathcal{U}_\alpha$. Then $\mathcal{B}^*$ and $\mathcal{B}^{*F}$ are open cover of $X$ clearly. $D_\alpha \subseteq \mathcal{U}_\alpha$ for each $\alpha$ and $X - \mathcal{V}_S = \bigcup_{\alpha \notin S} \mathcal{U}_\alpha$. $\mathcal{B}^*$ is an UP open cover and $\mathcal{B}^{*F}$ is an UP directed open cover of $X$. We call $\mathcal{B}^{*F}$ an **induced cover** of $\mathcal{B}$.
}


Let $\mathcal{B}$ be a LF family in $X$ and $\mathcal{B}^*$ be the induced cover as in Definition 1.2.2.

1. **(i)** If $\mathcal{B}^*$ has a cushioned open (co-cushioned) refinement, then $\mathcal{B}$ has a LF (point-finite) open expansion.


2. **(ii)** If $\mathcal{B}^*$ has a $\sigma$-cushioned open ($\sigma$-co-cushioned) refinement, then $\mathcal{B}$ has a $\sigma$-LF sequence ($\sigma$-point-finite sequence) of open expansions.



%#### Proof

We prove only one case; the other cases are...






$\mathcal{G} = \bigcup_{n<\omega} \mathcal{G}_{n}$ where each $\mathcal{G}_{n}$ is a cushioned open refinement of $\mathcal{U}$.

For each $n<\omega$, let $\mathcal{G}_{n} = \{G_{n,s} : s \in \Gamma\}$.

For each $\gamma < \tau$, $\Gamma \in \eta^{<\omega}$, $\overline{\bigcup_{s\in\Gamma} G_{n,s}} \subset \bigcup_{s\in\Gamma} U_{s}$. For $n<\omega$, $d\in D$, let

$$H_{n,d} = X \setminus \bigcup \{G_{n,s} : s\in \Gamma, d \notin s\}$$
Then $G_{n,s} \subset H_{n,d}$ for each $s \in \Gamma$ and $n<\omega$. For each $n<\omega$, $\mathcal{H}_{n} = \{H_{n,d} : d \in D\}$ is an open expansion of $D$. To see that $(\mathcal{H}_{n})_{n<\omega}$ is a sequence, let $x \in X$. There is $m<\omega$ and $s_{0} \in \Gamma$ such that $x \in G_{m,s_{0}}$. Then $x \notin H_{m,d}$ for all $d \notin s_{0}$. Thus $(\mathcal{H}_{n})_{n<\omega}$ is a sequence.


\end{document}
